\section{Custom block-sparse attention kernels}
\label{app:kernels}

The main text states only that we implement a custom efficient attention kernel;
this appendix documents its mechanism for a systems reviewer. Our cross-document
masks (\Cref{sec:method}) are highly structured but match none of the fixed sparsity
patterns supported by off-the-shelf kernels, so we implement a block-sparse
attention kernel in Triton \citep{tillet2019triton}. The design builds on
FlexAttention \citep{flexattention_blog2024,flexattention_paper2024} and the
FlashAttention line of tiled, IO-aware attention kernels
\citep{dao2022flashattention,dao2023flashattention2}, on the online-softmax
recurrence \citep{milakov2018onlinesoftmax} and memory-efficient attention
\citep{rabe2021selfattention} that underlie them, and on block-sparse GPU matmul
techniques \citep{gray2017blocksparse}. The distinguishing feature relative to
FlexAttention is that the sparse block schedule is precomputed \emph{once per batch}
and shared across all heads and layers, and the inner loop is specialized per block
class, rather than re-evaluating a mask predicate at every block during the kernel.
The cross-document-link mask uses the block-sparse kernel described below; the pure
document-causal mask uses a specialized variant of it (\S\ref{par:varlen}), and both
are selected automatically by mask type.

\paragraph{Block Interaction Mask (BIM).}
The central data structure is a precomputed, block-level compressed-sparse-row (CSR)
table recording, for each (query-block, key/value-block) pair, whether the two blocks
interact at all. It is analogous to FlexAttention's block mask but carries \emph{no
batch or head dimensions}: the attention mask is a pure function of the packed
sequence layout (document boundaries and cross-document grants), identical across all
heads and all layers. It is therefore built once per batch on the CPU (\char`~1--2\,ms
of vectorized NumPy) and reused by every head of every layer. The launch grid is
(\#blocks, $H$); each CTA walks only its own CSR row, so block pairs that never
interact are never launched and no per-block mask predicate is re-evaluated inside the
kernel. The table records the fraction of causal block pairs it is able to skip.

\paragraph{Bit-packed grant encoding.}
Cross-document grants (the rectangles of allowed cross-document attention induced by
detected links) are represented as bit-packed masks rather than as a dense $T\times T$
predicate. Grant $k$ is stored in bit $k \bmod 64$ of chunk $\lfloor k/64\rfloor$ of an
\texttt{int64} (bit 63 is written as \texttt{INT64\_MIN}, whose bit pattern is
correct under the \texttt{!=0} test); we maintain per-chunk query and key/value
bitmask tensors of shape $[\texttt{n\_chunks}, T]$, with
$\texttt{n\_chunks}=\lceil \texttt{max\_grants}/64\rceil$. Grant membership is then a
pointwise test with no reduction over the sequence: a query--key pair $(i,j)$ lies in
some grant iff
\[
\texttt{in\_grant}(i,j) \;=\; \bigvee_{c}\;\bigl(\texttt{q\_bm}[c,i]\ \texttt{\&}\ \texttt{kv\_bm}[c,j]\bigr)\neq 0,
\]
an AND per chunk followed by an OR across chunks (emitted as an unrolled loop). This
costs $O(T\cdot\texttt{n\_chunks})$ rather than $O(T^2)$: at $T=32\text{k}$ with the
production cap of $256$ grants ($\texttt{n\_chunks}=4$) the encoding is $2$\,MB, versus
roughly a gigabyte for a dense membership map. The same bitmasks are OR-reduced to
block granularity when building the BIM, letting the CSR construction prune whole block
pairs by the block-level union.

\paragraph{Block-class taxonomy and inner-loop specialization.}
Interacting block pairs differ in how much per-element masking they require, so the CSR
is partitioned into four classes, each with a specialized inner loop, in increasing
order of masking cost:
\begin{itemize}
  \item \textbf{Full} --- an off-diagonal pair of single-document blocks in the same
  document: no per-element mask, a dense tile matmul.
  \item \textbf{Pure-cross} --- single-document blocks in \emph{different} documents
  interacting only through a grant: same-document is false for every position pair by
  construction, so only the bitmask membership test is applied and the document-id
  tensor is never loaded.
  \item \textbf{Boundary} --- a block straddling a document boundary: the full
  per-element predicate $(\texttt{same\_doc}\lor\texttt{in\_grant})$ is evaluated.
  \item \textbf{Diagonal} --- an on-diagonal block: the full predicate plus the causal
  (lower-triangular) mask.
\end{itemize}
The CSR row is emitted as \texttt{[full, pure\_cross, boundary, diagonal]} so each CTA
processes progressively more expensive blocks in contiguous runs, keeping the branch
predictable; the transposed key/value$\to$query orientation used by the $dK/dV$
backward leads with the diagonal instead. The diagonal entry is forced to interact
even when a block's only valid positions are layout-gap tokens, since a query block
with no interactions would otherwise trigger an out-of-bounds CSR read in the backward.

\paragraph{Numerical correctness (two invariants).}
The block-level \texttt{same\_doc} test is an interval-overlap check on per-position
document labels, which is exact only when those labels are monotonic; because
graph-traversal packing makes the raw document ids non-monotonic, positions are
relabeled by an ordinal \emph{run index} $r(i)=\sum_{j\le i}\mathbb{1}[d_j\neq d_{j-1}]$
(the cumulative count of document-id changes), which is monotonic by construction and
equal for two positions iff they lie in the same contiguous run --- without it a
straddling query block silently drops its same-document keys, collapsing the softmax
support and diverging training (observed NaN loss with $dQ\approx 5.7\times 10^{4}$ on
document-linked code training). Separately, a query row with zero valid key/value
entries keeps the flash recurrence's initial sentinel maximum ($M=-10^{6}$), so the
backward $\exp_2(\text{score}-\text{LSE})$ overflows bf16 to $\infty$ and the ensuing
$\infty\times 0$ poisons the gradient; since the correct output of an empty row is
zero, a post-kernel \texttt{nan\_to\_num} maps such rows back to zero. Both are genuine
correctness fixes, not performance tweaks, and the second is a reusable
numerical-stability point for flash-style attention under full row masking.

\paragraph{Asymmetric forward/backward block sizes and low-level layout.}
The forward pass uses $128$-token blocks and the backward pass uses $64$-token blocks:
at head dimension $D_h=128$ a $128$-token backward tile exceeds the A100's
shared-memory and register budget, and the smaller backward tile is what keeps the
$D_h=128$ backward within budget. Two further choices reduce memory traffic. The
matmuls run natively in bfloat16 on the tensor cores with no up-cast to fp32 on load
and the softmax scale applied \emph{after} the product, halving the register and
shared-memory footprint of the operand tiles and raising occupancy. And the kernel
loads directly from the packed token--head--dimension layout using its native strides,
avoiding the permute/contiguous copies a $[B,H,T,D]$ layout would require; this saves
on the order of $200$\,MB of HBM traffic per step at $T=32\text{k}$ and brings peak
activation memory to FlexAttention parity.

\paragraph{Varlen document-causal specialization.}\label{par:varlen}
For the pure document-causal mask (no cross-document grants) a dedicated variant reuses
the same BIM CSR machinery. With no grants, all grant bitmasks are identically zero, so
the inner kernels skip \emph{all} bitmask loads and check only \texttt{same\_doc} and
causality: off-diagonal same-document blocks take the full fast path and only the
diagonal blocks are masked. The BIM is cached across steps, and a separate entry point
builds it directly from a per-position document-id vector, handling non-contiguous
document spans and layout-gap positions ($\texttt{doc\_id}=-1$). This gives the
throughput of a fused varlen self-attention forward while retaining reliable gradients
on real packed data.

\paragraph{Speed (as measured on A100).}
\Cref{tab:kernel-speed} reports throughput exactly as recorded in the repository's
kernel benchmark; all figures are measured on a single A100 80GB in bfloat16 against
compiled FlexAttention (\texttt{torch.compile(flex\_attention, dynamic=True)}) on a
representative $16$-head, $D_h=64$, document-length-$512$ configuration, and we
reproduce rather than re-benchmark them here. The forward pass is about $2\times$
faster ($1.14$ vs $2.38$\,ms at $T=32\text{k}$) and the backward about $12\%$ faster
($5.52$ vs $6.30$\,ms), for roughly $0.77\times$ total runtime at $T=32\text{k}$ (and
$0.91\times$ at $T=4\text{k}$). The most dramatic difference appears at head dimension
$D_h=128$, where FlexAttention's backward collapses to $\sim 177$\,ms while ours stays
at $\sim 16.85$\,ms --- a $9$--$10\times$ speedup, and the regime that motivates the
$64$-token backward tile above. Because the deployed cross-document kernel differs from
the benchmarked configuration only by the numerical guards described here (which do not
change the matmul work), these figures are representative of the deployed kernel but
were not separately re-measured for it. The tie-in to the density-aware batch scheduler
(the per-batch non-empty-block count) is described in \Cref{app:density} and not
repeated here.

\begin{table}[t]
  \centering
  \small
  \begin{tabular}{@{}llrrr@{}}
    \toprule
    Config & Pass & Ours & FlexAttention & Ratio \\
    \midrule
    $D_h{=}64$, $T{=}32\text{k}$ & forward  & $1.14$\,ms  & $2.38$\,ms  & $\sim2\times$ \\
    $D_h{=}64$, $T{=}32\text{k}$ & backward & $5.52$\,ms  & $6.30$\,ms  & $\sim1.12\times$ \\
    $D_h{=}64$, $T{=}32\text{k}$ & total    & ---         & ---         & $0.77\times$ \\
    $D_h{=}64$, $T{=}4\text{k}$  & total    & ---         & ---         & $0.91\times$ \\
    $D_h{=}128$                  & backward & $16.85$\,ms & $177$\,ms   & $\sim9$--$10\times$ \\
    \bottomrule
  \end{tabular}
  \caption{Kernel throughput, reproduced as-measured from the repository's benchmark
  (not re-run for this paper). All measurements are on a single A100 80GB in bfloat16;
  ratios are ours relative to compiled FlexAttention (lower total ratio is faster).
  ``total'' ratios are reported directly by the benchmark without per-pass timings.}
  \label{tab:kernel-speed}
\end{table}
